%%%%%%%%%%%%%%%%%%%%%%%%%%%%%%%%%%%%%%%%%%%%%%%%%%%%%%%%%%%%%%%%%%%%%%%%%%%%%%%%%
% Waterloo Rocketry Standard template for operations procedures                 %
% Copyright claimed by Aaron Morrison (akmorris@uwaterloo.ca)                   %
% Any commercial entity wishing to use this template in any way must provide    %
% to Waterloo Rocketry, one two-four (24 pack) of a non-shitty                  %
% type of beer or an amount of Canadian dollars roughly equalling               %
% the current cost of such a purchase. If you do so, then use it                %
% however you want, and I claim in no way that it compiles properly,            %
% and I take no responsibilty for any loss you incurr from using                %
% this template, and make no promises to support you in any technical           %
% capacity for sustained use.                                                   %
%%%%%%%%%%%%%%%%%%%%%%%%%%%%%%%%%%%%%%%%%%%%%%%%%%%%%%%%%%%%%%%%%%%%%%%%%%%%%%%%%
%%%%%%%%%%%%%%%%%%%%%%%%%%%%%%%%%%%%%%%%%%%%%%%%%%%%%%%%%%%%%%%%%%%%%%%%%%%%%%%%%
% Waterloo Rocketry Standard template for operations procedures                 %
% Copyright claimed by Aaron Morrison (akmorris@uwaterloo.ca)                   %
% Any commercial entity wishing to use this template in any way must provide    %
% to Waterloo Rocketry, one two-four (24 pack) of a non-shitty                  %
% type of beer or an amount of Canadian dollars roughly equalling               %
% the current cost of such a purchase. If you do so, then use it                %
% however you want, and I claim in no way that it compiles properly,            %
% and I take no responsibilty for any loss you incurr from using                %
% this template, and make no promises to support you in any technical           %
% capacity for sustained use.                                                   %
%%%%%%%%%%%%%%%%%%%%%%%%%%%%%%%%%%%%%%%%%%%%%%%%%%%%%%%%%%%%%%%%%%%%%%%%%%%%%%%%%
\documentclass[letter]{article}

%% margins and fonts
\usepackage[margin=1in]{geometry}
\renewcommand{\familydefault}{\sfdefault}

% begin package imports -- {{{
\usepackage{amsmath}
\usepackage{graphicx}
\usepackage[dvipsnames,table]{xcolor}
\usepackage{tabto}
\usepackage{titling}
\usepackage[iso,german]{isodate}
\usepackage{tabu}
% end package imports -- }}}

% begin set section title styles -- {{{
\usepackage{titlesec}
\titleformat{\section} {\setcounter{checklistnum}{0} \normalfont\Large\bfseries}{}{0em}{}[{\titlerule[0.8pt]}]
\titleformat{\subsection} {\setcounter{checklistnum}{0} \normalfont\large}{}{0em}{}[{\titlerule[0.6pt]}]
% end set section title styles -- }}}

% begin checklist symbols definition -- {{{
\usepackage{enumitem,amssymb}
\newcounter{checklistnum}
\setcounter{checklistnum}{0}
\DeclareRobustCommand{\checklistnumber}{\refstepcounter{checklistnum}\thechecklistnum}
\newlist{checklist}{itemize}{6}
\setlist[checklist,1]{
label={\color{gray}\checklistnumber}\hspace{2em}$\square$,
leftmargin=0em,
itemindent=2em
}
\setlist[checklist,2]{
label={\color{gray}\checklistnumber}\hspace{4em}$\square$,
leftmargin=0em,
itemindent=4em
}
\setlist[checklist,3]{
label={\color{gray}\checklistnumber}\hspace{6em}$\square$,
leftmargin=0em,
itemindent=6em
}
\setlist[checklist,4]{
label={\color{gray}\checklistnumber}\hspace{8em}$\square$,
leftmargin=0em,
itemindent=8em
}
\setlist[checklist,5]{
label={\color{gray}\checklistnumber}\hspace{10em}$\square$,
leftmargin=0em,
itemindent=10em
}
\setlist[checklist,6]{
label={\color{gray}\checklistnumber}\hspace{12em}$\square$,
leftmargin=0em,
itemindent=12em
}
% end checklist symbols definition -- }}}

% begin personnel macro -- {{{
\newcommand{\operator}[4]{%
  \expandafter\newcommand\csname #1\endcsname{{\color{#2}\textbf{#3}}\phantom{}}
  \expandafter\newcommand\csname #1full\endcsname{{\color{#2}\textbf{#4 [#3]}}\phantom{}}
}
% end personnel macro -- }}}

% begin valve macro -- {{{
\newcommand{\valve}[3]{%
  \expandafter\newcommand\csname #1\endcsname{#2 (#3)}
  \expandafter\newcommand\csname #1num\endcsname{#3}
}
% end valve macro --  }}}

% begin flag macro -- {{{
\newcommand\flag[1]{
    \hfill \makebox(0,0){\hspace{1.7in}\textcolor{#1}{\rule{0.3in}{0.5in}}}
}
% end flag macro -- }}}

\pagenumbering{gobble}


\title{
\huge Eridium Liquid Rocket Engine\\ Static Fire 2\\ 
\vspace{1cm}
\large Static Fire Operations Procedure}



\begin{document}



% begin titlepage -- {{{
\begin{center}
\vspace*{7cm}
\hspace{7em}\includegraphics[width=30em]{common/mono_horizontal_standard}
\newline
\rule{50em}{2pt}

\vspace{1cm}
\thetitle

\vspace*{\fill}
Compiled on \today
\end{center}
\newpage
% end titlepage -- }}}


\section{Static Fire Test Operations Procedure}

% begin operators declarations -- {{{
% syntax for operator macro:
% arg1: their name (so here you use \auth to insert the AUTHOR name)
% arg2: the color you want them highlighted in
% arg3: their abbreviated title (what you should use in the checklists)
% arg4: their full title (insert into document with arg1 + full, so for auth
%       it'd be "\authfull" to insert their full title). Use this in Personnel
%       Required subsection, and in Sign Off subsection.
\operator{ops}{YellowOrange}{OPS}{Operations Director}
\operator{daq}{Cerulean}{DAQ}{Data Acquisition System Operator}
\operator{rlcs}{blue}{CONTROL}{Control System Operator}
\operator{primary}{ForestGreen}{PRIMARY}{Primary Operator}
\operator{secondary}{violet}{SECONDARY}{Secondary Operator}
\operator{emergency}{red}{EMERGENCY}{Emergency Services Caller}
\operator{peri}{red}{P1}{Perimeter Guard 1}
\operator{perii}{red}{P2}{Perimeter Guard 2}
\operator{periii}{red}{P3}{Perimeter Guard 3}
\operator{periv}{red}{P4}{Perimeter Guard 4}
% end operators declarations -- }}}

% begin valve declaration -- {{{
% syntax for valve macro:
% arg1: Abbreviated name (to save typing out the full callout every time)
% arg2: Valve description
% arg3: Valve P&ID tag
\valve{oxfill}{Remote Oxidizer Fill Valve}{OV-101}
\valve{oxdrain}{Remote Oxidizer Drain Valve}{OV-102}
\valve{vent}{Oxidizer Tank Vent Valve}{OV-202}
\valve{purge}{Nitrogen Purge Valve}{NV-101}
\valve{moxfill}{Manual Oxidizer Fill Bypass Valve}{OMV-101}
\valve{moxdrain}{Manual Oxidizer Drain Valve}{OMV-102}
\valve{fiso}{Manual Fuel Tank Isolation Valve}{FMV-102}
\valve{pdrain}{Manual Pneumatics Drain Valve}{NMV-101}
\valve{fdrain}{Manual Fuel Dump Valve}{FMV-101}
\valve{oxinj}{Oxidizer Injector Valve}{OV-301}
\valve{finj}{Fuel Injector Valve}{FV-101}
\valve{inj}{Injector Valve}{IJV-101}
\valve{oxiso}{Oxidizer Tank Isolation Valve}{OV-201}
\valve{ohshit}{Manual Oxidizer Tank Vent Valve}{OMV-201}
% end valve declaration -- }}}

% begin contents subsection -- {{{
\subsection{Contents}
This document contains the following four procedures:
\begin{itemize}
    \item The \textit{Fill System Leak Check Procedure} comprises steps for verifying the integrity of the system plumbing and proper operation of the data acquisition and test control systems, using nitrogen gas.
    \item The \textit{Fill System Check Abort} comprises steps for aborting and safing the system in the event of an anomaly during the Fill System Check Procedure
    \item The \textit{Static Fire Test Procedure} comprises steps for operating the fill system and conducting a static fire test of the engine.
    \item The \textit{Static Fire Abort} comprises steps for aborting and safing the system in the event of an anomaly during the Static Fire Test procedure. 
\end{itemize}
% end contents subsection -- }}}

% begin personnel required section -- {{{
\subsection{Personnel Required}
The test operations team consists of ten personnel:
\begin{itemize}
    \item The \opsfull{} directs operations and communicates with all other test personnel
    \item The \primaryfull{} is the main system operator. \primary{} operates all manual valves on the fill and nitrogen pressurization system.
    \item The \secondaryfull{} is the backup for \primary{}, and communicates with \ops{}. If \primary{} becomes incapacitated, \secondary{} is responsible for
        removing them from danger.
    \item The \rlcsfull{} operates the test control system, including actuating remote valves and engine ignition.
    \item The \daqfull{} monitors and operates the data acquisition system.
    \item \perifull{}, \periifull{}, \periiifull{} and \perivfull{} ensure no unauthorized personnel enter the testing area during test operations.
    \item The \emergencyfull{} is responsible for calling 911 in the event of an emergency. 
\end{itemize}
\setcounter{checklistnum}{0}
% end personnel required section -- }}}


\subsection{General Notes}
\begin{itemize}
\item In the event of loss of radio contact, personnel should contact \ops{} by phone at 249-385-6237. 
\item For \primary{} and \secondary{}, if contact with \ops{} cannot be re-established in 1 minute, close all cylinders and return to mission control.
\item In the event of an emergency requiring 911 to be called, first responders should be directed to \ops{} before entering the site for a briefing on hazards.  
\item Unless otherwise specified, any event failing to occur is grounds for abort. 
\end{itemize}

\newpage
% begin sign off -- {{{
\subsection{Sign-Off}
\textit{To be completed by all test personnel after reading and familiarization with procedures}
\begin{itemize}
    \item \opsfull{}            \tabto{20em}\rule{20em}{0.4pt}
    \item \primaryfull{}        \tabto{20em}\rule{20em}{0.4pt}
    \item \secondaryfull{}      \tabto{20em}\rule{20em}{0.4pt}
    \item \rlcsfull{}           \tabto{20em}\rule{20em}{0.4pt}
    \item \daqfull{}            \tabto{20em}\rule{20em}{0.4pt}
    \item \perifull{}           \tabto{20em}\rule{20em}{0.4pt}
    \item \periifull{}          \tabto{20em}\rule{20em}{0.4pt}
    \item \periiifull{}         \tabto{20em}\rule{20em}{0.4pt}
    \item \perivfull{}          \tabto{20em}\rule{20em}{0.4pt} 
    \item \emergencyfull{}      \tabto{20em}\rule{20em}{0.4pt}
\end{itemize}
\setcounter{checklistnum}{0}
% end sign off -- }}}

% begin prior to start section -- {{{
\newpage
\subsection{Prior to Fill System Leak Check Start}
\begin{checklist}
    \item Ensure the following items are complete before beginning the test procedure
    \begin{checklist}
        \item Injector assembly.
        \item Concentric tank assembly, including pre-filling fuel tank with 3.75kg of ethanol.
        \item Plumbing setup per Eridium LSF2 Leak Check P\&ID.
        \item Data acquisition system setup.
        \item Valve safe state checks.
        \item RLCS power on, including:
        \begin{checklist}
        \item Both lipos installed.
        \item Actuator and main switches flipped.
        \item GSPD on and set to AC mode.
        \end{checklist}
        \item Perimeter sweep. 
        \item Primary and secondary vent line checks.
        \item Correct plumbing is connected to cylinders and caps are in place.
        \item Generator filled with fuel and running. 
        \item System control key inserted in towerside box.
        \item Set all regulators to the minimum position. 
        \item Ropes installed on all manual ball valves and tested. 
        \item Place an appropriately sized container below \fdrain{} to collect ethanol \\ \phantom{} \hspace{8ex} in the event dumping is needed.
        \item Low pressure leak check using nitrogen at approx. 100psi and soapy water. 
        \item Ensure compressor is turned on and fully pressurized.
        \item Ensure a spray bottle of soapy water is at the test stand.
        \item Verify all valve ropes are fully unrolled. 
        \item Ensure all personnel as defined above are available and have completed sign-off.
        \item Ensure a blank copy of the procedure is available or a unique symbol will be used to cross off steps. 
        \item Verify the lockout key is inserted into the towerside box.
        \item Ensure the following personnel have walkie-talkies and that communication is functional:
        \begin{checklist}
            \item \ops{}
            \item \secondary{}
            \item \peri{}
            \item \perii{}
            \item \periii{}
            \item \periv{}
    \end{checklist}
    \item Ensure \peri{}, \perii{}, \periii{}, and \periv{} are wearing hi-vis vests.
    \item Ensure \primary{} and \secondary{} are wearing face shields and have no exposed skin.
    \item Ensure \primary{} is wearing thermal gloves.
    \item Remove propane cylinder from shipping container, if present.
    \end{checklist}
    
\end{checklist}
\setcounter{checklistnum}{0}
% end prior to start section -- }}}

\subsection{Fill System Leak Check Procedure}
\begin{checklist}
    \item \primary{}: Open \fdrain{}. 
    \item \primary{}: Fill ethanol through \fdrain{} until no more can be added. Be careful not to spill any ethanol. 
    \item \primary{}: Close \fdrain{}.
    \item \primary{}: Open \fiso{}.
    \item \daq{}: Zero load cell. 
    \item \daq{}: Begin Data logging. 
    \item \daq{}: Confirm the zero points of all sensors have been recorded. 
    \item \primary{}: Confirm the following valves are initially closed:
    \begin{checklist}
        \item \oxfill{}. 
        \item \moxfill{}.
        \item \oxdrain{}.
        \item \moxdrain{}. 
        \item \ohshit{}.
        \item \purge{}.
        \item \fdrain{}.
        \item \pdrain{}.
    \end{checklist}
    \item \primary{}: Confirm the following valves are initially open:
    \begin{checklist}
        \item \fiso{}.
        \item \oxiso{}.
    \end{checklist}
    \item \rlcs{}: Confirm the commanded state for the following valves is 'closed':
    \begin{checklist}
        \item \inj{}.
    \end{checklist}
    \item \rlcs{}: Confirm the commanded state for the following valves is 'open':
    \begin{checklist}
        \item \vent{}.
    \end{checklist}
    \item \daq{}: If state feedback is available, confirm the following valves are initially closed:
    \begin{checklist}
        \item \finj{}.
        \item \oxinj{}.
    \end{checklist}
    \item \secondary{}: Confirm there are no personnel in the test exclusion zone except for \primary{} and \\ \phantom{} \hspace{3ex} \secondary{}.
    \item \peri{}, \perii{}, \periii{}, \periv{}: Close the perimeter and do not allow any individuals to enter the test exclusion zone.
    \item \primary{}: Remove covers from plumbing in the following locations:
    \begin{checklist}
        \item Oxidizer Fill Panel PRV (OPRV-101).
        \item Oxidizer fill panel drain hose.
        \item Oxidizer run tank PRV (OPRV-201).
        \item Nitrogen Fill panel PRV (NPRV-101).
        \item \vent{}.
        \item \fdrain{}.
        \item \pdrain{}
    \end{checklist}
    \item \primary{}: Slowly adjust compressor regulator NREG-201 until a pressure of approximately 145psi is read \\ \phantom{} \hspace{3ex} by the downstream gauge (NPG-202). 
    \item \primary{}: Communicate final pneumatics pressure as read by NPG-201.
    \item \daq{}: Confirm pressure reading on NPT-101 is in agreement. 
    \item \primary{}: Slowly open nitrogen supply cylinder NSC-101.
    \item \primary{}: Comunicate cylinder pressure as read by NPG-101.
    \item \primary{}: Slowly adjust pneumatics regulator NREG-101 until a pressure of 750psi is read on the \\ \phantom{} \hspace{3ex} downstream gauge (NPG-102).
    \item \daq{}: Confirm pressure reading on OPT-101 is in agreement. 
    \item \primary{} and \secondary{}: Spray soapy water on all fittings.
    \item \secondary{}: Remove lockout key from towerside box.
    \item \primary{} and \secondary{}: Retreat to safe distance. 
    \item \rlcs{}: Insert system control key into clientside box and engage control.
    \item \rlcs{}: Close \vent{}.
    \begin{checklist}[label=$\bullet$]
        \item If \ventnum{} fails to close, close \oxiso.
        \item If \oxisonum{} fails to close, proceed to the Fill System Leak Check Abort Procedure.
    \end{checklist}
    \item \rlcs{}: Open \oxfill{}. 
    \begin{checklist}[label=$\bullet$]
        \item If \oxfillnum{} fails to open proceed to the Fill System Leak Check Abort Procedure.
    \end{checklist}
    \item \daq{}: Communicate fuel and oxidizer tank pressure continuously as read by OPT-201 and FPT-101. 
    \item \ops{}: Proceed when oxidizer tank pressure reaches pressure of 750psi.
    \item \ops{}: Record pressure difference between OPT-201 and FPT-101. 
    \item \rlcs{}: Close \oxfill{}.
    \item \daq{}: Monitor oxidizer and fuel tank pressures. 
     \begin{checklist}[label=$\bullet$]
        \item If leaks are observed, close \oxiso and determine if the leak rate slows. 
        \item If leak rate does not slow after 5 minutes, proceed.
    \end{checklist}
    \item \rlcs{}: Open \vent{}.
    \item \rlcs{}: Open \oxiso{} if closed. 
    \item \rlcs{}: Open \oxdrain{}.
    \item \ops{}: Proceed when oxidizer tank is fully vented. 
    \item \rlcs{}: Disengage control and remove system control key from clientside box. 
    \item \primary{} and \secondary{}: Return to the test stand. 
    \item \secondary{}: Insert lockout key into towerside box. 
    \item \primary{}: Inspect all plumbing for evidence of bubbles. Report findings to \ops{}.
    \item \ops{}: If leaks were discovered that can be easily fixed, have \primary{} tighten the fittings and return to \\ \phantom{} \hspace{3ex} Line 39. 
    \item \primary{}: Close Nitrogen supply cylinder NSC-101.
    \item \primary{}: Set compressor regulator NREG-201 to minimum position. 
    \item \primary{}: Slowly open \pdrain{}.
    \item \secondary{}: Remove lockout key from towerside box. 
    \item \primary{} and \secondary{}: Retreat to safe distance. 
    \item \rlcs{}: Insert system control key into clientside box and engage control. 
    \item \rlcs{}: Open \oxfill{}.
    \item \daq{}: Confirm all sensors are reading their determined zero points. 
    \item \rlcs{}: Disengage control and remove system control key from the clientside box.
    \item \peri{}, \perii{}, \periii{}, and \periv{}: Open the perimeter. 
    \item \daq{}: Stop DAQ logging and save data files. 
    \item \primary{} and \secondary{}: Approach the test stand.  
    \item \primary{}: Replace covers on plumbing in the following 7 locations:
    \begin{checklist}
        \item Oxidizer fill panel PRV (OPRV-101).
        \item Oxidizer fill panel drain hose.
        \item Oxidizer run tank PRV (OPRV-201).
        \item Nitrogen Fill panel PRV (NPRV-101)
        \item \vent{}.
        \item \fdrain{}.
        \item \pdrain{}
    \end{checklist}    
    \item \ops{}: If leaks were observed an not fixed, correct leaks and repeat procedure as needed. Once no leaks are \\ \phantom{} \hspace{3ex} present, proceed with final setup for Static Fire Test. 
\end{checklist}
\setcounter{checklistnum}{0}


\newpage
\subsection{Fill System Leak Check Abort Procedure}
\begin{checklist}
 \item \rlcs{}: Open \vent{} if not already open.
 \begin{checklist}[label=$\bullet$]
    \item If RLCS is not operational, have \primary{} open \moxdrain{} using the \\ \phantom{} \hspace{8ex} attached rope. 
    \end{checklist}
 \item \rlcs{}: Open \oxiso{} if not already open. 
 \item \rlcs{}: Close \oxfill{} if not already closed. 
 \item \rlcs{}: Open \oxdrain{}. 
 \item \rlcs{}: Remove system control key and disengage control.
 \item \ops{}: Proceed when oxidizer tank is fully vented. 
 \item \primary{} and \secondary: Approach the test stand. 
 \item \secondary{}: Insert lockout key into towerside box. 
 \item \primary{}: Close nitrogen supply cylinder NSC-101. 
 \item \primary{}: Turn off the air compressor.
 \item \primary{}: Set the air compressor regulator to zero. 
 \item \secondary{}: Remove lockout key from towerside box.
 \item \primary{} and \secondary{}: Retreat to a safe distance. 
 \item \rlcs{}: Insert system control key and engage control. 
 \item \rlcs{}: Open \oxfill{}. 
 \begin{checklist}[label=$\bullet$]
    \item If RLCS is not operational, have \primary{} open \moxfill{} using \\ \phantom{} \hspace{8ex} the attached rope. 
    \end{checklist}
\item \rlcs{}: Open \purge{}. 
\begin{checklist}[label=$\bullet$]
    \item If RLCS is not operational, proceed. Nitrogen will be drained through \\ \phantom{} \hspace{8ex} \pdrain{}. 
    \end{checklist}
\item \rlcs{}: Disengage control and remove the system control key. 
\item \daq{}: Confirm all sensors are reading their determined zero points. 
\begin{checklist}[label=$\bullet$]
    \item If DAQ is not operational, have \primary{} verify the gauges OPG-101, OPG-201 and NPG-101 are all \\ \phantom{} \hspace{8ex} reading zero. 
    \end{checklist}
\item \primary{}: Open \pdrain{}.
\item \peri{}, \perii{}, \periii{}, and \periv{}: Open the perimeter.
\item \ops{}: Abort and revisit the plumbing and controls setup. 
\end{checklist}


\newpage
% begin first checklist-- {{{ 
\subsection{Prior to Static Fire Test Start}
% syntax for the flag macro:
% arg1: The colour you want the flag to be
% this should be useful for abort procedures that we need to put sticky notes on for irec judges to happy for
% It prints the flag all the way over the right edge so the flag should be visible even in the middle of a stack of ops
\begin{checklist}
\item Ensure the following items are complete before beginning the test procedure
    \begin{checklist}
        \item Fill System Leak Check procedure complete.
        \item Combustion chamber assembly.
        \item Plumbing setup per Eridium LSF2 P\&ID. 
        \item Data acquisition system setup.
        \item Valve safe state checks.
        \item RLCS power on, including:
        \begin{checklist}
        \item Both lipos installed.
        \item Actuator and main switches flipped.
        \item GSPD on and set to AC mode.
        \end{checklist}
        \item Perimeter sweep. 
        \item Primary and secondary vent line checks.
        \item Correct plumbing is connected to cylinders and caps are in place.
        \item Tank heating jackets installed on oxidizer run tank and set to 60\textdegree C. 
        \item Cameras in the correct positions and prepared to record. 
        \item Generator filled with fuel and running. 
        \item System control key inserted in towerside box.
        \item Set all regulators to the minimum position. 
        \item Ensure ropes are installed on manual ball valves and tested.
        \item Place an appropriately sized container below \fdrain{} to collect ethanol \\ \phantom{} \hspace{8ex} in the event dumping is needed.
        \item Ensure fuel tank is filled with 3.75kg of ethanol. 
        \item Ensure all personnel as defined above are available and have completed sign-off.
         \item Ensure a blank copy of the procedure is available or a unique symbol will be used to cross off steps. 
         \item Verify lockout key is installed in towerside box. 
         \item Verify ignition cables are disconnected from the towerside box.
         \item Verify a multimeter is at the test tand.
        \item Verify all valve ropes are fully unrolled. 
        \item Ensure the following personnel have walkie-talkies and that communication is functional:
        \begin{checklist}
            \item \ops{}
            \item \secondary{}
            \item \peri{}
            \item \perii{}
            \item \periii{}
            \item \periv{}
        \end{checklist}
        \item Ensure \peri{}, \perii{}, \periii{}, and \periv{} are wearing hi-vis vests.
        \item Ensure \primary{} and \secondary{} are wearing face shields and have no exposed skin.
        \item Ensure \primary{} is wearing thermal gloves.
        \item Remove propane heater from shipping container, if present.
    \end{checklist}
      
\end{checklist}
\setcounter{checklistnum}{0}
% end first checklist -- }}}

\subsection{Static Fire Test Procedure}
\begin{checklist}
    \item \daq{}: Begin DAQ data logging. 
    \item \secondary{}: Confirm that the ignition wire is not connected to the towerside control box.
    \item \daq{}: Confirm the zero points of all sensors have been recorded. 
    \item \primary{}: Confirm the following valves are initially closed:
    \begin{checklist}
        \item \oxfill{}. 
        \item \moxfill{}.
        \item \oxdrain{}.
        \item \moxdrain{}.
        \item \purge{}. 
        \item \pdrain{}.
        \item \ohshit{}.
        \item \fdrain{}.
    \end{checklist}
    \item \primary{}: Confirm the following valves are initially open:
    \begin{checklist}
        \item \fiso{}.
        \item \oxiso{}.
    \end{checklist}
    \item \rlcs{}: Confirm the commanded state for the following valves is 'closed':
    \begin{checklist}
        \item \inj{}.
    \end{checklist}
    \item \rlcs{}: Confirm the commanded state for the following valves is 'open':
    \begin{checklist}
        \item \vent{}. 
    \end{checklist}
    \item \daq{}: If state feedback is available, confirm the following valves are initially closed:
    \begin{checklist}
        \item \finj{}. 
        \item \oxinj{}.
    \end{checklist}
    \item \primary{}: Confirm the resistance across the primary and secondary ignition cables is 2.0 +/- 0.5 ohms. 
    \item \secondary{}: Confirm there are no personnel in the test exclusion zone except for \primary{} and \\ \phantom{} \hspace{3ex} \secondary{}.
    \item \daq{}: Confirm all sensors are reading their determined zero points. 
    \item \peri{}, \perii{}, \periii{}, \periv{}: Close the perimeter and do not allow any individuals to enter the test exclusion zone.
    \item \primary{}: Confirm all covers are removed from the system. There should be 8 covers total, with covers \\ \phantom{} \hspace{3ex} in the following locations:
    \begin{checklist}
        \item Oxidizer fill panel PRV (OPRV-101).
        \item Oxidizer fill panel drain hose.
        \item Oxidizer run tank PRV (OPRV-201).
        \item Nitrogen Fill Panel PRV (NPRV-101).
        \item \vent{}.
        \item \pdrain{}.
        \item \fdrain{}.
        \item Combustion chamber/nozzle.
    \end{checklist}
    \item \primary{}: Slowly open Nitrogen Supply Cylinder NSC-101. 
    \item \primary{}: Communicate cylinder supply pressure as read by NPG-101. 
    \item \primary{}: Slowly adjust pneumatics regulator NREG-101 until a pressure of 145psi is read on NPG-102.
    \item \daq{}: Confirm the pressure reading on NPT-101 is 145psi. 
    \item \primary{}: Slowly open the first Nitrous Oxide supply cylinder, OSC-101. 
     \begin{checklist}[label=$\bullet$]
        \item If leaks are observed proceed to the Static Fire Remote Abort Procedure. 
    \end{checklist}
    \item \primary{}: Slowly open the second Nitrous Oxide supply cylinder, OSC-102. 
    \item \primary{}: Communicate the oxidizer supply pressure as read by gauge OPG-101.
    \item \daq{}: Communicate the oxidizer supply pressure as read by OPT-101.
    \item \ops{}: Confirm the two values are in agreement. 
    \item \primary{}: Start recording on all manual cameras. 
    \item \secondary{}: Connect the ignition wire to the towerside box. 
    \item \secondary{}: Remove lockout key from towerside box. 
    \item \primary{} and \secondary{}: Retreat to mission control. Bring the multimeter.  
    \item \rlcs{}: Insert system control key into clientside box and engage control.
    \item \rlcs{}: Open \oxdrain{}.
    \item \rlcs{}: Open \purge{} to begin nitrogen purge of plumbing.
    \item \ops{}: Proceed after 1 minute.
    \item \rlcs{}: Close \oxdrain{}.
    \item \rlcs{}: Close \purge{}.
    \item \rlcs{}: Open \oxfill{}.
    \begin{checklist}[label=$\bullet$]
        \item If \oxfill{} fails to the Static Fire Remote Abort Procedure.
    \end{checklist}
    \item \daq{}: Continuously communicate oxidizer tank mass as read by oxidizer tank load cell. 
    \item \ops{}: Proceed when tank mass has increased by 13.5 kg (+ 0.5kg, - 0 kg), or a plume is observed from  \\ \phantom{} \hspace{3ex} \vent{}. 
    \begin{checklist}[label=$\bullet$]
        \item If supply and oxidizer tank pressure equilibrate before target conditions are met proceed to the Static \\ \phantom{} \hspace{8ex} Fire Remote Abort Procedure.
    \end{checklist}
    \item \rlcs{}: Close \oxfill{}.
    \begin{checklist}[label=$\bullet$]
        \item If \oxfill{} fails to close proceed to the Static Fire Remote Abort Procedure.
    \end{checklist}
    \item \rlcs{}: Close \vent{}.
    \begin{checklist}[label=$\bullet$]
        \item If \ventnum{} fails to close, close \oxiso.
        \item If \oxiso{} fails to close, proceed to the Static Fire Remote Abort \\ \phantom{} \hspace{8ex} Procedure.
    \end{checklist}
    \item \rlcs{}: Engage tank heating.
    \item \daq{}: Communicate oxidizer tank pressure as read by OPT-201 in increments of 50psi.
    \item \daq{}: Monitor all sensors and report any deviations or anomalies to \ops{} immediately. 
    \item \ops{}: If oxidizer tank mass is decreasing or leaking from \vent{} is suspected, \\ \phantom{} \hspace{3ex} close \oxiso{}. 
    \item \ops{}: If oxidizer tank mass drops below 12 kg proceed to the Static Fire abort procedure,
    \item \ops{}: Proceed when the oxidizer tank pressure has reached 700psi. 
    \item \ops{}: Ensure all spectators are wearing safety glasses and have earplugs in. 
    \item \ops{}: Verify all remote cameras are recording. 
    \item \ops{}: Instruct drone operator to launch. 
    \item \ops{}: Proceed when the oxidizer tank pressure has reached firing pressure of 750psi.  
    \item \rlcs{}: Disengage tank heating. 
    \item \textbf{PAUSE POINT}
    \item \ops{}: Verify all test commit criteria are met. 
    \item \textbf{BEGIN RADIO SILENCE} - Note that in the event of an emergency or anomaly requiring immediate \\ \phantom{} \hspace{3ex} attention radio silence will be broken. 
    \item \rlcs{}: Conduct the engine startup procedure. 
        \begin{checklist}
         \item Arm primary ignition.
         \item  Hold down the fire button and continuously communicate ignition current. 
         \item Proceed when either primary or secondary ignition current reads 0 amps, or when smoke is observed.
         \begin{checklist}[label=$\bullet$]
            \item In the event of a failed ignition (criteria  are not met within 1 minute) immediately release the \\ \phantom{} \hspace{12ex} ignition button and proceed to the Static Fire Remote Abort \\ \phantom{} \hspace{12ex} Procedure
            \end{checklist}
         \item Conduct the static fire test by opening \inj{}. 
            \begin{checklist}[label=$\bullet$]
            \item In the event \inj{} does not open proceed to the Static Fire Remote Abort \\ \phantom{} \hspace{12ex} Procedure.
            \end{checklist}
         \item \textbf{ALL}: Observe the plume.
         \end{checklist}
    \item \ops{}: Proceed when the static fire has concluded.
    \item \daq{}: Verify DAQ systems are functional and all sensors are transmitting data. 
    \item \ops{}: Confirm with remote camera operators that there is no visible fire on or around the test stand. 
    \begin{checklist}[label=$\bullet$]
            \item In the event fire is observed:
             \begin{checklist}
                \item \ops{}: If the fire is major or outside of the fill area, \textbf{call 911 immediately}.
                \item \ops{}:  If fire is localized to the engine, proceed with nitrogen purge. Repeat as needed until the \\ \phantom{} \hspace{12ex} fire is extinguished. 
                \item If a small fire is present in the fill area that cannot be extinguishers by a nitrogen purge, send  \\ \phantom{} \hspace{12ex} \primary{} and \secondary{} to use fire extinguishers. 
                \end{checklist} 
            \end{checklist}
    \item \rlcs{}: Conduct a nitrogen purge by opening \purge{}. 
    \item \rlcs{}: Close \purge{}. 
    \item \rlcs{}: Close \inj{}. 
    \item \textbf{END RADIO SILENCE}
    \item \rlcs{}: Open \vent{}. 
    \item \rlcs{}: Open \oxiso{} if closed. 
    \item \rlcs{}: Disengage control and remove system control key.
    \item \daq{}: Confirm CC temperature as read by CCTC-101, -102, -103, and/or -104 is below 100\textdegree C, unless \\ \phantom{} \hspace{3ex} suspected faulty. 
    \item \primary{} and \secondary{}: Approach the test stand. 
    \item \secondary{}: Insert lockout key into towerside box. 
    \item \primary{}: Close Nitrous Oxide supply cylinder OSC-101.
    \item \primary{}: Close Nitrous Oxide supply cylinder OSC-102. 
    \item \primary{}: Close Nitrogen supply cylinder NSC-101.
    \item \secondary{}: Remove lockout key from towerside box.
    \item \primary{} and \secondary{}: Retreat to mission control. 
    \item \rlcs{}: Insert system control key into clientside box and engage control. 
    \item \rlcs{}: Open \oxdrain{}{}.
    \item \rlcs{}: Open \oxfill{}{}.
    \item \rlcs{}: Open \purge{}. 
    \item \daq{}: Confirm all sensors are reading their determined zero points. 
    \item \rlcs{}: Disengage control and remove system control key from the clientside box.
    \item \daq{}: Stop DAQ logging and save data files. 
    \item \primary{} and \secondary{}: Approach the test stand. 
    \item \secondary{}: Insert lockout key into towerside box. 
    \item \primary{}: Open \pdrain{}.
    \item \peri{}, \perii{}, \periii{}, and \periv{}: Open the perimeter. 
    \item \primary{}: Replace plumbing covers in the following 7 locations:
    \begin{checklist}
        \item Oxidizer fill panel PRV (OPRV-101).
        \item Oxidizer fill panel drain hose.
        \item Oxidizer run tank PRV (OPRV-201).
        \item Nitrogen Panel PRV (NPRV-101)
        \item \vent{}.
        \item \fdrain{}.
        \item \pdrain{}.
    \end{checklist}
    \item \daq{}: Continue to monitor CC temperature. Inform \ops{} immediately if temperature exceeds 100\textdegree C. 
    \item \ops{}: Proceed with teardown. 
\end{checklist}
    
\newpage
\subsection{Static Fire Remote Abort Procedure}
\begin{checklist}
   \item \rlcs{}: Release the ignition button if pressed.
    \item \daq{}: Verify DAQ is operational and receiving data. 
    \begin{checklist}[label=$\bullet$]
            \item If DAQ or RLCS is not operational, proceed to the Static Fire Manual Abort Procedure. 
            \end{checklist}
    \item \rlcs{}: Disengage tank heating. 
    \item \rlcs{}: Open \vent{} if not already open. 
    \item \rlcs{}: Open \oxiso{} if not already open. 
    \item \rlcs{}: Close \oxfill{} if not already closed. 
    \item \rlcs{}: Close \purge{} if not already closed. 
    \item \rlcs{}: Open \oxdrain{}.
    \item \ops{}: Proceed when the oxidizer tank is fully vented. 
    \item \rlcs{}: Close \inj{} if not already closed.
    \item \rlcs{}: Disengage control and remove the system control key. 
    \item \primary{} and \secondary{}: Approach the test stand. 
    \item \secondary{}: Insert lockout key into towerside box. 
    \item \secondary{}: Disconnect the ignition wires from the towerside box. 
    \item \primary{}: Close Nitrous Oxide supply cylinder OSC-101. 
    \item \primary{}: Close Nitrous Oxide supply cylinder OSC-102. 
    \item \primary{}: Close Nitrogen supply cylinder NSC-101. 
    \item \secondary{}: Remove lockout key from towerside box.
    \item \primary{} and \secondary{}: Retreat to a safe distance. 
    \item \rlcs{}: Insert the system control key and engage control. 
    \item \rlcs{}: Open \oxfill{}.
    \item \rlcs{}: Open \purge{}. 
    \item \daq{}: Verify all sensors are reading their determined zero point. 
    \item \rlcs{}: Disengage control and remove system control key. 
    \item \primary{} and \secondary{}: Approach the test stand. 
    \item \primary{}: Open \pdrain{}. 
    \begin{checklist}[label=$\bullet$]
            \item If ethanol draining is required, have \primary{} place a bucket underneath the outlet hose from \\ \phantom{} \hspace{7ex} \fdrain{}, then open \fdrain{}. 
    \end{checklist}
    \item \peri{}, \perii{}, \periii{}, and \periv{}: Open the perimeter. 
    \item \ops{}: Abort and revisit plumbing and controls setup.
\end{checklist}

\newpage
\subsection{Static Fire Manual Abort Procedure}
\begin{checklist}
     \item \rlcs{} and \daq{}: Attempt to troubleshoot the issue. 
      \begin{checklist}[label=$\bullet$]
            \item If RLCS and DAQ become operational, have \primary{} and \secondary{} retreat and return to the \\ \phantom{} \hspace{7ex} Static Fire Remote Abort Procedure. 
            \item If a connection issue is suspected send a DAQ technician with a laptop and walkie talkie to the manual \\ \phantom{} \hspace{7ex} valve pull ropes and attempt to connect to the local wifi. If successful, proceed with a remote abort. 
    \end{checklist}
    \item \primary{} and \secondary{}: Approach the manual valve ropes and listen for venting. 
    \begin{checklist}[label=$\bullet$]
            \item If venting is heard, report approximate source to \ops{}, and do not proceed until cleared \\ \phantom{} \hspace{7ex} by \ops{}.
    \end{checklist}
    \item \primary{}: Open \moxdrain{} using the attached rope. 
    \begin{checklist}[label=$\bullet$]
            \item If no venting is heard or \moxdrain{} fails to open, open \\ \phantom{} \hspace{7ex} \ohshit{}. 
    \end{checklist}
    \item \ops{}: Proceed when \primary{} and \secondary{} report no venting is heard.
    \item \secondary{}: Insert lockout key into towerside box. 
    \item \secondary{}: Disconnect the ignition wires from the towerside box. 
    \item \primary{}: Close Nitrous Oxide supply cylinder OSC-101. 
    \item \primary{}: Close Nitrous Oxide supply cylinder OSC-102. 
    \item \primary{}: Close Nitrogen supply cylinder NSC-101. 
    \item \primary{} and \secondary{}: Retreat to the manual valve ropes. 
    \item \primary{}: Open \moxfill{}.  
    \item \ops{}: Proceed when \primary{} and \secondary{} report no venting is heard.
    \item \primary{}: Open \pdrain{}. 
    \item \primary{}: Verify the following gauges read zero pressure:
    \begin{checklist}[label=$\bullet$]
        \item Oxidizer fill panel gauge OPG-101.
        \item Nitrogen fill panel gauge NPG-101. 
        \item Vent panel gauge OPG-201. 
    \end{checklist}
    \begin{checklist}[label=$\bullet$]
            \item If ethanol draining is required, have \primary{} place a bucket underneath the outlet hose from \\ \phantom{} \hspace{7ex} \fdrain{}, then open \fdrain{}. 
    \end{checklist}
    \item \peri{}, \perii{}, \periii{}, and \periv{}: Open the perimeter. 
    \item \ops{}: Abort and revisit plumbing and controls setup. 
\end{checklist}


\newpage
\subsection{Static Fire Emergency Abort Procedure}
\begin{checklist}
    \item \rlcs{}: Release the ignition button if pressed. 
    \item \rlcs{}: Open \oxdrain{}. 
    \begin{checklist}[label=$\bullet$]
            \item If DAQ or RLCS is not operational and it is safe to approach, have \primary{} open \\ \phantom{} \hspace{7ex} \moxdrain{} using the attached rope. 
            \item If no venting occurs when \\ \phantom{} \hspace{7ex} \moxdrain{} is opened or \moxdrain{} fails to \\ \phantom{} \hspace{7ex} open, have \primary{} open \ohshit{} using the attached rope.
            \item If is it not safe to approach the test site, call 911. 911 should be directed to \ops{} before approaching \\ \phantom{} \hspace{7ex} the site. 
    \end{checklist}
    \item \rlcs{}: Close \oxfill{} if not already closed. 
    \item \rlcs{}: Open \vent{} if not already open. 
    \item \rlcs{}: Open \oxiso{} if not already open.
    \item \rlcs{}: Disengage tank heating.
    \begin{checklist}[label=$\bullet$]
            \item If DAQ or RLCS is not operational, have \primary{} turn off the generator.
    \end{checklist}
    \item \ops{}: Proceed when the oxidizer run tank is fully vented.
    \item \rlcs{}: Close \inj{} if not already closed.
    \item \rlcs{}: Remove system control key and disengage control. 
    \item \primary{} and \secondary{}: Approach the test stand. 
    \item \secondary{}: Insert the lockout key into the towerside box.
    \item \secondary{}: Disconnect the ignition wires from the towerside box. 
    \item \primary{}: Close Nitrous Oxide supply cylinder OSC-101. 
    \item \primary{}: Close Nitrous Oxide supply cylinder OSC-102. 
    \item \primary{}: Close Nitrogen supply cylinder NSC-101. 
    \item \secondary{}: Remove system control key from the towerside box.
    \item \primary{} and \secondary{}: Retreat to a safe distance. 
    \item \rlcs{}: Insert the system control key and engage control. 
    \item \rlcs{}: Open \oxfill{}.
    \item \rlcs{}: Open \purge{}. 
    \item \daq{}: Verify all sensors are reading their determined zero point.
        \begin{checklist}[label=$\bullet$]
            \item If DAQ is not operational,\primary{} and \secondary{}. are not to approach the test stand if venting \\ \phantom{} \hspace{7ex} is heard.
    \end{checklist}
    \item \rlcs{}: Disengage control and remove system control key. 
    \item \primary{} and \secondary{}: Approach the test stand. 
    \item \primary{}: Open \pdrain{}. 
    \begin{checklist}[label=$\bullet$]
            \item If ethanol draining is required, have \primary{} place a bucket underneath the outlet hose from \\ \phantom{} \hspace{7ex} \fdrain{}, then open \fdrain{}. 
    \end{checklist}
    \item \peri{}, \perii{}, \periii{}, and \periv{}: Open the perimeter. 
    \item \ops{}: Abort and revisit plumbing and controls setup. 
\end{checklist}

\newpage
% begin Variable Recording section -- {{{
\subsection{Variable Recording}
\renewcommand{\arraystretch}{2}
\begin{tabular}{ | m{15em} | m{15em}| m{15em} |} 
    \hline
    PT1 (Ox Fill) zero point &  & \\
    \hline
    PT2 (Ox Tank) zero point & &  \\
    \hline
    PT3 (Nitrogen Fill) zero point & &  \\
    \hline
    PT4 (Fuel Tank) zero point & &   \\
    \hline
    PT5 (Pneumatics) zero point & &   \\
    \hline
    PT6 (Ox injector) zero point &  &  \\
    \hline
    PT7 (Fuel injector) zero point & & \\
    \hline
    PT8 (CC) zero point & & \\
    \hline
    PT9 (CC) zero point & &\\
    \hline
    Ox tank load cell zero point &  & \\
    \hline
    Fuel tank load cell zero point & &\\
    \hline
    CC load cell zero point &  & \\
    \hline
    Pneumatic pressure &  & \\
    \hline
     &  & \\
    \hline
& &  \\
    \hline

       &  & \\
    \hline
       &  & \\
    \hline
       &   & \\
    \hline
       & &  \\
    \hline
       & &  \\
    \hline
       & &  \\
    \hline
       & &  \\
    \hline
       &  & \\
    \hline
       &  & \\
    \hline
       &  & \\
    \hline
       & &  \\
    \hline
       &  & \\
    \hline

\end{tabular}
% end Variable Recording section }}}


\end{document}
